A trustworthy explanation of a dashcam collision must name the right vehicles, state the right cause, and place the impact at the right moment.
We release \dataset{}, a dataset of 1{,}500 dashcam clips for which our annotators wrote the complete forensic ground truth---severity, impact geometry, vehicles by colour and type, weather, crash window, and a multi-sentence causal explanation---and use it to separate two ways Video-LLMs fail: \emph{omission} (leaving out what happened) and \emph{hallucination} (asserting what did not).
Four open Video-LLMs evaluated zero-shot on a 150-video held-out split are fluent but omit 28--98\% of forensic fields.
QLoRA adaptation of Qwen2.5-VL-7B on our 1{,}198 training videos (\crashlogic) raises BLEU-4 from $0.016$ to $0.142$ and cuts the omission cost from $0.462$ to $0.107$ ($p<10^{-4}$).
The hallucination cost does not move ($0.227\to0.223$, $p=0.81$): one in five adapted explanations names a vehicle whose colour contradicts the ground truth, and the model fabricates a collision in every confirmed no-crash test video.
Its apparently strong crash-window localisation (tIoU $0.012\to0.373$) is a learned dataset prior: only two windows are ever emitted, the predicted start is uncorrelated with the true start, and a constant window scores higher.
Finally, Temporal Contrastive Decoding, which subtracts logits from a reversed-frame negative, does not repair the residue: mild contrast changes no metric significantly and strong contrast significantly increases both failure modes.
Reversal preserves every order-invariant cue, so the contrast cannot reach the dominant errors.
We release the annotations, outputs, per-video scores, and an error taxonomy.
